\chapter{Trưởng Thành Là Bớt Đổ Lỗi, Tăng Chịu Trách Nhiệm}
\markboth{Trưởng Thành Là Bớt Đổ Lỗi, Tăng Chịu Trách Nhiệm}{Maturity Means Less Blame and More Responsibility}

\section{Trưởng Thành Là Bớt Đổ Lỗi}
\begin{truyen}{Trưởng Thành Là Bớt Đổ Lỗi}{Growing Up Means Blaming Less}
\chuhoa{N}{gười trẻ có rất nhiều lý do thật để khó khăn: gia cảnh, nền tảng, môi trường, và cả bất công.}
\textit{(Young people have many real reasons for difficulty: family background, limited support, environment, and unfairness.)}

Nhưng nếu mọi chuyện trong đời đều được giải thích chỉ bằng lỗi của người khác, em sẽ mất dần quyền thay đổi phần của mình.
\textit{(But if everything in life is explained only by other people’s fault, you slowly lose the power to change your own part.)}

Trưởng thành không phải phủ nhận khó khăn thật.
\textit{(Growing up does not mean denying real hardship.)}

Nó là biết trong hoàn cảnh đó, phần nào là phần mình vẫn phải chịu trách nhiệm.
\textit{(It means knowing what part remains your responsibility within that hardship.)}

\end{truyen}

\begin{baihoc}
Đổ lỗi đôi khi giúp em đỡ đau trong chốc lát. Chịu trách nhiệm mới giúp em đổi đời về lâu dài.
\textit{(Blame may reduce pain for a moment. Responsibility is what changes life over the long run.)}
\end{baihoc}

\begin{ghinhoanh}
\textbf{Short English to remember:}

Blame may reduce pain for a moment.

Responsibility is what changes life over the long run.

\end{ghinhoanh}

\begin{tuhocnhanh}
\textbf{Cau hoi tu hoc / Self-study questions}

1. Van de chinh la gi? \textit{What is the main problem?}

2. Bai hoc thuc te la gi? \textit{What is the practical lesson?}

3. Mot cau tieng Anh em muon nho la cau nao? \textit{Which English sentence do you want to remember?}
\end{tuhocnhanh}
\ngancach

\section{Xin Lỗi Và Sửa Lỗi Là Hai Việc Khác Nhau}
\begin{truyen}{Xin Lỗi Và Sửa Lỗi Là Hai Việc Khác Nhau}{Apologizing and Repairing Are Two Different Things}
\chuhoa{C}{ó người xin lỗi rất hay nhưng vẫn lặp lại cùng một lỗi mãi.}
\textit{(Some people apologize well but keep repeating the same mistake.)}

Lời xin lỗi khi đó dần mất giá.
\textit{(Then the apology slowly loses value.)}

Người trưởng thành không chỉ biết nhận sai bằng miệng.
\textit{(A mature person does not only admit fault with words.)}

Họ sửa sai bằng cách đổi hành vi.
\textit{(They repair it by changing behavior.)}

Nói đúng mà sống không đổi thì vẫn chưa đủ.
\textit{(Speaking correctly without changing life is still not enough.)}

\end{truyen}

\begin{baihoc}
Muốn lớn thật, đừng chỉ học nói câu xin lỗi. Hãy học cả cách làm người khác thấy mình đã khác đi.
\textit{(To truly mature, do not only learn the words “I am sorry.” Learn how to let people see that you have changed.)}
\end{baihoc}

\begin{ghinhoanh}
\textbf{Short English to remember:}

To truly mature, do not only learn the words “I am sorry.” Learn how to let people see that you have changed.

\end{ghinhoanh}

\begin{tuhocnhanh}
\textbf{Cau hoi tu hoc / Self-study questions}

1. Van de chinh la gi? \textit{What is the main problem?}

2. Bai hoc thuc te la gi? \textit{What is the practical lesson?}

3. Mot cau tieng Anh em muon nho la cau nao? \textit{Which English sentence do you want to remember?}
\end{tuhocnhanh}
\ngancach

\section{Dọn Phần Bừa Bộn Mình Tạo Ra}
\begin{truyen}{Dọn Phần Bừa Bộn Mình Tạo Ra}{Clean Up the Mess You Create}
\chuhoa{M}{ột dấu hiệu của trưởng thành là không bỏ lại hậu quả của mình cho người khác dọn mãi.}
\textit{(One sign of maturity is not leaving your consequences for other people to clean forever.)}

Đi trễ thì xin lỗi và sửa lịch.
\textit{(If you are late, apologize and fix your schedule.)}

Làm hỏng việc thì nhận phần mình.
\textit{(If you damage work, own your part.)}

Tiêu quá tay thì tự siết lại.
\textit{(If you overspend, tighten your own budget.)}

Người lớn không được định nghĩa bằng tuổi.
\textit{(Adults are not defined by age.)}

Họ được định nghĩa bằng khả năng dọn phần bừa bộn mình gây ra.
\textit{(They are defined by their ability to clean up the mess they create.)}

Ai cứ muốn hưởng quyền mà né phần dọn dẹp thì mới chỉ lớn cái xác.
\textit{(Those who want rights but avoid cleanup have grown only in body.)}

\end{truyen}

\begin{baihoc}
Chịu trách nhiệm không hào nhoáng, nhưng nó làm con người có trọng lượng hơn rất nhiều.
\textit{(Responsibility is not glamorous, but it gives a person much more weight and substance.)}
\end{baihoc}

\begin{ghinhoanh}
\textbf{Short English to remember:}

Responsibility is not glamorous, but it gives a person much more weight and substance.

\end{ghinhoanh}

\begin{tuhocnhanh}
\textbf{Cau hoi tu hoc / Self-study questions}

1. Van de chinh la gi? \textit{What is the main problem?}

2. Bai hoc thuc te la gi? \textit{What is the practical lesson?}

3. Mot cau tieng Anh em muon nho la cau nao? \textit{Which English sentence do you want to remember?}
\end{tuhocnhanh}
\ngancach

\section{Không Còn Tuổi Để Chờ Ai Nhắc Mãi}
\begin{truyen}{Không Còn Tuổi Để Chờ Ai Nhắc Mãi}{There Comes an Age When No One Should Keep Reminding You}
\chuhoa{L}{ớn lên rồi, có những việc em không thể chờ người khác nhắc mãi mới làm.}
\textit{(As you grow older, there are tasks you cannot keep waiting to be reminded about.)}

Đúng giờ, giữ lời, trả tiền, phản hồi, và chuẩn bị là những thứ dần phải thành phản xạ.
\textit{(Punctuality, keeping promises, paying what you owe, replying, and preparing must gradually become reflexes.)}

Nếu cái gì cũng cần ai đó nhắc, em vẫn đang sống như người khác cầm tay mình đi.
\textit{(If everything still requires reminders, you are living as if someone else must hold your hand.)}

Một phần trưởng thành là biến điều đúng thành thói quen, không phải biến nó thành bài giảng người khác phải lặp đi lặp lại.
\textit{(Part of maturity is turning what is right into habit, not into a lecture others must repeat forever.)}

\end{truyen}

\begin{baihoc}
Tự nhắc được mình là bước rất quan trọng của trưởng thành. Người nào cũng chờ bị nhắc thì rất khó lớn thật.
\textit{(Being able to remind yourself is a major step in growing up. People who always wait to be reminded rarely mature deeply.)}
\end{baihoc}

\begin{ghinhoanh}
\textbf{Short English to remember:}

Being able to remind yourself is a major step in growing up.

People who always wait to be reminded rarely mature deeply.

\end{ghinhoanh}

\begin{tuhocnhanh}
\textbf{Cau hoi tu hoc / Self-study questions}

1. Van de chinh la gi? \textit{What is the main problem?}

2. Bai hoc thuc te la gi? \textit{What is the practical lesson?}

3. Mot cau tieng Anh em muon nho la cau nao? \textit{Which English sentence do you want to remember?}
\end{tuhocnhanh}
\ngancach

\section{Sống Làm Sao Để Sau Này Khỏi Xấu Hổ Với Mình}
\begin{truyen}{Sống Làm Sao Để Sau Này Khỏi Xấu Hổ Với Mình}{Live So That Future You Will Not Be Ashamed}
\chuhoa{C}{uối cùng, trưởng thành không chỉ là kiếm được tiền hay sống riêng được.}
\textit{(In the end, maturity is not only about earning money or living on your own.)}

Nó còn là chuyện mỗi ngày em đang thành loại người nào.
\textit{(It is also about what kind of person you are becoming each day.)}

Có người lớn lên về tuổi nhưng nhỏ đi về nhân cách.
\textit{(Some people grow older in age while shrinking in character.)}

Có người đi chậm, còn vụng, còn nghèo, nhưng mỗi năm lại đáng tin hơn và đàng hoàng hơn.
\textit{(Some move slowly, remain clumsy, remain poor, yet each year become more trustworthy and decent.)}

\end{truyen}

\begin{baihoc}
Đích của trưởng thành không phải chỉ là sống sót. Nó là sống sao để không phải cúi mặt trước chính mình.
\textit{(The goal of adulthood is not only survival. It is living in a way that does not make you lower your eyes before yourself.)}
\end{baihoc}

\begin{ghinhoanh}
\textbf{Short English to remember:}

The goal of adulthood is not only survival.

It is living in a way that does not make you lower your eyes before yourself.

\end{ghinhoanh}

\begin{tuhocnhanh}
\textbf{Cau hoi tu hoc / Self-study questions}

1. Van de chinh la gi? \textit{What is the main problem?}

2. Bai hoc thuc te la gi? \textit{What is the practical lesson?}

3. Mot cau tieng Anh em muon nho la cau nao? \textit{Which English sentence do you want to remember?}
\end{tuhocnhanh}
\ngancach

\section{Trả Lại Điều Mình Đã Mượn}
\begin{truyen}{Trả Lại Điều Mình Đã Mượn}{Return What You Borrowed}
\chuhoa{M}{ột dấu hiệu rất thật của trưởng thành là không xem nợ nhỏ hay lời hứa nhỏ là chuyện không đáng nhớ.}
\textit{(A very real sign of maturity is not treating small debts or small promises as forgettable.)}

Nhiều người trẻ tưởng thiếu một chút rồi xin khất thêm cũng không sao.
\textit{(Many young people think a little delay on borrowed things does not matter much.)}

Nhưng cách em xử lý những món nhỏ thường cho thấy em coi trọng trách nhiệm đến đâu.
\textit{(But how you handle small obligations reveals how seriously you treat responsibility.)}

Người đáng tin được nhận ra từ những việc bé như vậy.
\textit{(Trustworthy people are recognized through such small acts.)}

\end{truyen}

\begin{baihoc}
Trưởng thành bắt đầu từ lúc em không để người khác phải đòi phần em đã nhận.
\textit{(Maturity begins when others do not need to chase what you owe.)}
\end{baihoc}

\begin{ghinhoanh}
\textbf{Short English to remember:}

Maturity begins when others do not need to chase what you owe.

\end{ghinhoanh}

\begin{tuhocnhanh}
\textbf{Cau hoi tu hoc / Self-study questions}

1. Van de chinh la gi? \textit{What is the main problem?}

2. Bai hoc thuc te la gi? \textit{What is the practical lesson?}

3. Mot cau tieng Anh em muon nho la cau nao? \textit{Which English sentence do you want to remember?}
\end{tuhocnhanh}
\ngancach

\section{Giữ Lời Với Việc Nhỏ}
\begin{truyen}{Giữ Lời Với Việc Nhỏ}{Keep Your Word in Small Things}
\chuhoa{C}{ó người giữ lời hứa lớn rất cảm động nhưng lại quên những việc rất nhỏ mình từng nói sẽ làm.}
\textit{(Some people speak beautifully about big promises but forget very small things they said they would do.)}

Họ tưởng chỉ việc lớn mới đáng tính là trách nhiệm.
\textit{(They think only large matters count as responsibility.)}

Nhưng cuộc sống chung được giữ bằng vô số lời hứa nhỏ được hoàn thành đều đặn.
\textit{(But shared life is held together by many small promises completed steadily.)}

Ai xem nhẹ chuyện nhỏ thường cũng làm người khác mệt dần trong chuyện lớn.
\textit{(Those who dismiss small duties often tire others out in bigger matters too.)}

\end{truyen}

\begin{baihoc}
Người lớn đáng tin là người làm xong cả phần nhỏ không ai vỗ tay.
\textit{(A reliable adult completes even the small parts that earn no applause.)}
\end{baihoc}

\begin{ghinhoanh}
\textbf{Short English to remember:}

A reliable adult completes even the small parts that earn no applause.

\end{ghinhoanh}

\begin{tuhocnhanh}
\textbf{Cau hoi tu hoc / Self-study questions}

1. Van de chinh la gi? \textit{What is the main problem?}

2. Bai hoc thuc te la gi? \textit{What is the practical lesson?}

3. Mot cau tieng Anh em muon nho la cau nao? \textit{Which English sentence do you want to remember?}
\end{tuhocnhanh}
\ngancach

\section{Nhận Ra Giới Hạn Của Mình}
\begin{truyen}{Nhận Ra Giới Hạn Của Mình}{Admitting Your Own Limits}
\chuhoa{C}{ó lúc trưởng thành không phải là gồng lên nhận hết, mà là nói rõ mình chưa gánh nổi.}
\textit{(Sometimes maturity is not taking everything on, but saying clearly that you cannot carry it yet.)}

Nhiều người trẻ tưởng nhận hết mới là bản lĩnh.
\textit{(Many young people think accepting everything is proof of strength.)}

Nhưng nhận quá sức rồi làm đổ mới là cách nhanh nhất để mất uy tín.
\textit{(But taking more than you can carry and dropping it is the fastest way to lose credibility.)}

Biết giới hạn giúp em phân phối sức đúng hơn và học nhanh hơn.
\textit{(Knowing your limits helps you use energy more wisely and learn faster.)}

\end{truyen}

\begin{baihoc}
Chịu trách nhiệm không có nghĩa ôm tất cả. Nó là biết ôm đúng phần mình làm được.
\textit{(Responsibility does not mean carrying everything. It means carrying your real part well.)}
\end{baihoc}

\begin{ghinhoanh}
\textbf{Short English to remember:}

Responsibility does not mean carrying everything.

It means carrying your real part well.

\end{ghinhoanh}

\begin{tuhocnhanh}
\textbf{Cau hoi tu hoc / Self-study questions}

1. Van de chinh la gi? \textit{What is the main problem?}

2. Bai hoc thuc te la gi? \textit{What is the practical lesson?}

3. Mot cau tieng Anh em muon nho la cau nao? \textit{Which English sentence do you want to remember?}
\end{tuhocnhanh}
\ngancach

\section{Đừng Chờ Người Lớn Đến Cứu Mãi}
\begin{truyen}{Đừng Chờ Người Lớn Đến Cứu Mãi}{Stop Waiting for Adults to Fix Everything}
\chuhoa{C}{ó một độ tuổi em phải thôi chờ người lớn hơn giải quyết giúp các rối rắm mình tạo ra.}
\textit{(There is an age when you must stop waiting for older people to fix the problems you created.)}

Trước đó nhiều người vẫn tưởng chỉ cần thú nhận rồi sẽ có ai đó đứng ra dọn tiếp.
\textit{(Before that, many still believe admitting the problem is enough and someone else will clean up the rest.)}

Nhưng trưởng thành bắt đầu khi em tự bước vào phần dọn hậu quả của mình.
\textit{(But maturity begins when you step into the cleanup yourself.)}

Người lớn thật không chỉ biết báo lỗi.
\textit{(Real adults do not only report mistakes.)}

Họ biết sửa lỗi.
\textit{(They repair them.)}

\end{truyen}

\begin{baihoc}
Tự chịu phần bừa bộn mình gây ra là một mốc lớn của việc thành người lớn.
\textit{(Cleaning the mess you create is a major milestone in becoming an adult.)}
\end{baihoc}

\begin{ghinhoanh}
\textbf{Short English to remember:}

Cleaning the mess you create is a major milestone in becoming an adult.

\end{ghinhoanh}

\begin{tuhocnhanh}
\textbf{Cau hoi tu hoc / Self-study questions}

1. Van de chinh la gi? \textit{What is the main problem?}

2. Bai hoc thuc te la gi? \textit{What is the practical lesson?}

3. Mot cau tieng Anh em muon nho la cau nao? \textit{Which English sentence do you want to remember?}
\end{tuhocnhanh}
\ngancach

\section{Sống Để Mai Sau Nhìn Lại Không Xấu Hổ}
\begin{truyen}{Sống Để Mai Sau Nhìn Lại Không Xấu Hổ}{Live So You Will Not Be Ashamed Later}
\chuhoa{C}{uối cùng, có những quyết định không ai chấm ngay nhưng vài năm sau chính em sẽ là người phải nhìn lại chúng.}
\textit{(In the end, some decisions are not judged right away, but years later you will look back at them yourself.)}

Nhiều người trẻ tưởng qua được hôm nay là đủ.
\textit{(Many young people think surviving today is enough.)}

Nhưng đời dài khiến mỗi lựa chọn dần trở thành nhân cách.
\textit{(But over a long life, each choice slowly becomes character.)}

Sống có trách nhiệm là sống sao để tương lai của mình không phải cúi mặt vì hiện tại của mình.
\textit{(Living responsibly means living so your future self does not lower its eyes before your present self.)}

\end{truyen}

\begin{baihoc}
Đích của trưởng thành không chỉ là sống sót. Nó là sống cho đàng hoàng.
\textit{(The goal of maturity is not only survival. It is decent living.)}
\end{baihoc}

\begin{ghinhoanh}
\textbf{Short English to remember:}

The goal of maturity is not only survival.

It is decent living.

\end{ghinhoanh}

\begin{tuhocnhanh}
\textbf{Cau hoi tu hoc / Self-study questions}

1. Van de chinh la gi? \textit{What is the main problem?}

2. Bai hoc thuc te la gi? \textit{What is the practical lesson?}

3. Mot cau tieng Anh em muon nho la cau nao? \textit{Which English sentence do you want to remember?}
\end{tuhocnhanh}
